\thispagestyle{empty}
\begin{tikzpicture}[remember picture,overlay]
  \fill[Ink] (current page.south west) rectangle (current page.north east);
  \draw[Accent,line width=1.2pt]
    ([xshift=18mm,yshift=-22mm]current page.north west)
    rectangle
    ([xshift=-18mm,yshift=22mm]current page.south east);
  \foreach \x/\y in {35/62,58/76,78/55,112/72,142/54,165/76}{
    \fill[Accent] ([xshift=\x mm,yshift=\y mm]current page.south west) circle (2.2pt);
  }
  \draw[Rule,opacity=.6,line width=.6pt]
    ([xshift=35mm,yshift=62mm]current page.south west) --
    ([xshift=58mm,yshift=76mm]current page.south west) --
    ([xshift=78mm,yshift=55mm]current page.south west) --
    ([xshift=112mm,yshift=72mm]current page.south west) --
    ([xshift=142mm,yshift=54mm]current page.south west) --
    ([xshift=165mm,yshift=76mm]current page.south west);
  \node[anchor=north west,text width=165mm]
    at ([xshift=24mm,yshift=-55mm]current page.north west) {
      \sffamily\color{white}\fontsize{35}{42}\selectfont\bfseries
      \BookTitle\\[8mm]
      \fontsize{16}{22}\selectfont\mdseries\color{Rule}
      \BookSubtitle
    };
  \node[anchor=south west]
    at ([xshift=24mm,yshift=34mm]current page.south west) {
      \sffamily\color{white}\large \BookEdition\quad\textcolor{Accent}{|}\quad\BookVersion
    };
\end{tikzpicture}
\null\clearpage

\thispagestyle{empty}
\vspace*{0.24\textheight}
{\sffamily\Huge\bfseries\color{Ink}\BookTitle\par}
\vspace{1em}
{\sffamily\Large\color{Muted}\BookSubtitle\par}
\vfill
{\sffamily\BookEdition\\\BookDate\par}
\clearpage

\thispagestyle{empty}
\vspace*{0.12\textheight}
{\sffamily\Large\bfseries About this edition\par}
\vspace{1em}
This mentoring edition is a guided textbook and laboratory manual. It teaches
the ideas beneath HTTP before asking the reader to build a server. The software
is intentionally small enough to inspect completely. It is not offered as a
production web server.

\vfill
\noindent \BookTitle\\
\BookSubtitle\\
\BookEdition, \BookVersion\\
\BookDate

\vspace{2em}
\noindent The examples use Python's standard library. Protocol descriptions
are based on the standards listed in the bibliography. Code may be used and
modified for educational purposes.
\clearpage

\chapter{How to use this book}

This book is both a roadmap and a construction manual. The first seven
chapters form a theory gate: do not rush through them simply because they do
not yet create a server. A parser is much easier to design when you can explain
what a byte stream is, where message boundaries come from, and which HTTP rules
exist to prevent two recipients from disagreeing about a request.

Beginning with Part II, every major step has a frozen source snapshot under
\filepath{code/stages}. Read a chapter, predict the behavior, run the snapshot,
make a small change, and then compare your result with the next stage. Part III
turns the reusable framework into a complete SQLite-backed Task API.

\begin{objectives}
By the end, you should be able to:
\begin{itemize}
  \item narrate a request from a URL through DNS, TCP, HTTP, routing, and storage;
  \item explain why TCP reads do not correspond to HTTP messages;
  \item parse a strict, length-framed HTTP/1.1 subset safely;
  \item design request, response, routing, middleware, and server abstractions;
  \item implement and test a resource-oriented JSON API;
  \item state exactly why the finished server is educational rather than production-ready.
\end{itemize}
\end{objectives}

\section{The weekly rhythm}

Plan for six to eight focused hours each week: two or three hours of reading,
two or three hours of code-along work, and one or two hours for exercises and
reflection. Keep a lab notebook. For every surprising result, write the bytes
you expected, the bytes you observed, and the rule that explains the difference.

\begin{center}
\begin{tabularx}{\textwidth}{@{}>{\bfseries}l X X@{}}
\toprule
Weeks & Focus & Deliverable \\
\midrule
1--4 & Network and HTTP theory & Annotated request journey and theory checkpoint \\
5--9 & Socket server and framework & Tested \code{tinyhttp} package \\
10--12 & SQLite REST API and hardening & Complete Task API and final assessment \\
\bottomrule
\end{tabularx}
\end{center}

\section{Conventions}

Wire messages use a monospaced block. The symbol \code{CRLF} means the two
bytes carriage return (13) and line feed (10), written in Python as
\code{b"\\r\\n"}. Commands use \code{python}; on Windows, \code{py} may be
the correct launcher. Paths use forward slashes, even though Windows also
accepts backslashes in many shell contexts.

\begin{securityboundary}
Run the examples on \code{127.0.0.1}. Binding to \code{0.0.0.0} exposes the
learning server to other machines that can reach your computer. Do not expose
it to the public Internet.
\end{securityboundary}

